% ══════════════════════════════════════════════════════════════════
%  NOTATION CONVENTION  –  final version for IJPR
%  Requires: \usepackage{mathabx} or custom \widecheck command
% ══════════════════════════════════════════════════════════════════

\begin{convention}[Index suppression and aggregation]
\label{conv:notation}

Variables in the model carry up to six indices (see
Table~\ref{tab:nomenclature}). To avoid writing three and four-level
summations and writing aditional constraints to imply boolena logic, we adopt the following convention: suppressing
one or more indices from a variable and placing an accent on its name
produces a derived variable whose values are obtained by
collapsing the suppressed indices according to one of three aggregation
rules. The accent identifies which rule is applied.

We call the original, fully-indexed variable the root variable.
The indices that remain after the suppression are called projected
indices, collected into the tuple~$\pi$. The set of all combinations of the suppressed indices is
denoted~$\mathcal{S}$, so that $|\mathcal{S}|$ is the number of root
entries sharing a given combination of retained indices.



\paragraph{Dot notation.}

%----------------------------------------------------------------------
\medskip
\noindent\textbf{Rule~1 \textnormal{(Sum,~$\widetilde{\phantom{x}}$).}}
Suppressing one or more indices and writing a tilde denotes the
sum over the suppressed indices:
\[
  \widetilde{\phi}_{\pi}
  \;\coloneqq\;
  \sum_{\sigma\in\mathcal{S}} \phi_{\pi,\sigma} \quad \forall \pi,
\]
for each combination of projected indices~$\pi$, where $\sigma$ ranges
over all $|\mathcal{S}|$ combinations of the suppressed indices. For
continuous and integer root variables the result is a total; for binary
root variables it is an integer count of active entries, not itself
binary.

\smallskip\noindent\textit{Example.}
From the root variable $NF_{ijmpr}$
$\bigl((i,j)\!\in\!E,\;m\!\in\!\mathcal{M},\;p\!\in\!P,
  \;r\!\in\!R\bigr)$, suppressing $p$
($\mathcal{S}=P$) and suppressing $p,r$
($\mathcal{S}=P\times R$):
\begin{align}
  \widetilde{NF}_{ijmr}
  &= \sum_{p\in P} NF_{ijmpr}
  && \forall\,(i,j)\in E,\;m\in\mathcal{M},\;r\in R,
  \label{eq:sum-ex1} \\[4pt]
  \widetilde{NF}_{ijm}
  &= \sum_{p\in P}\sum_{r\in R} NF_{ijmpr}
  && \forall\,(i,j)\in E,\;m\in\mathcal{M}.
  \label{eq:sum-ex2}
\end{align}
When a sum variable is introduced explicitly into the model, it is
linked to its root by the corresponding defining equality.

When two indices range over the same set, as $i,j\in V$ do in
$NF_{ijmpr}$, suppressing one by name alone is ambiguous.
A dot~$(\,\cdot\,)$ marks the suppressed position:
$\widetilde{NF}_{\cdot jmpr}\coloneqq\sum_{i\in V}NF_{ijmpr}$
suppresses~$i$, while
$\widetilde{NF}_{i\cdot mpr}\coloneqq\sum_{j\in V}NF_{ijmpr}$
suppresses~$j$. This convention applies to all three rules.


%----------------------------------------------------------------------
\medskip
\noindent\textbf{Rule~2
  \textnormal{(OR,~$\widecheck{\phantom{x}}$).}
  Binary root variables only.}
Suppressing one or more indices and writing a caron introduces a new
binary variable equal to~1 if and only if at least one root
binary sharing the projected indices is active. A new variable
$\widecheck{\phi}_{\pi}$ is coupled to the corresponding sum
$\widetilde{\phi}_{\pi}$ (Rule~1) via:
\begin{align}
  \widecheck{\phi}_{\pi}
    &\leq \widetilde{\phi}_{\pi},
    \label{eq:or-lb} \\
  \widetilde{\phi}_{\pi}
    &\leq |\mathcal{S}|\;\widecheck{\phi}_{\pi},
    \label{eq:or-ub}
\end{align}
for each combination of projected indices~$\pi$, where $|\mathcal{S}|$
is the number of combinations of the suppressed indices. Together,
\eqref{eq:or-lb}--\eqref{eq:or-ub} enforce
$\widecheck{\phi}_{\pi}=1 \iff \widetilde{\phi}_{\pi}\geq 1$.

\smallskip\noindent\textit{Example.}
From the binary $XF_{ijmpr}$, suppressing~$p$
($\pi=(i,j,m,r)$, $\mathcal{S}=P$, $|\mathcal{S}|=|P|$)
yields $\widecheck{XF}_{ijmr}$, which equals~1 whenever route
$(i,j)$ carries full RTIs by mode~$m$ and RTI type~$r$ for any
product:
\begin{align}
  \widecheck{XF}_{ijmr}
    &\leq \widetilde{XF}_{ijmr}
    && (i,j)\in E,\;m\in\mathcal{M},\;r\in R,
    \label{eq:or-ex-lb} \\
  \widetilde{XF}_{ijmr}
    &\leq |P|\;\widecheck{XF}_{ijmr}
    && (i,j)\in E,\;m\in\mathcal{M},\;r\in R,
    \label{eq:or-ex-ub}
\end{align}

%----------------------------------------------------------------------
\medskip
\noindent\textbf{Rule~3
  \textnormal{(At-most-one,~$\widehat{\phantom{x}}$).}
  Binary root variables only.}
Suppressing one or more indices and writing a circumflex introduces a
new binary variable that equals the sum and simultaneously caps it
at~1, asserting that at most one root binary sharing the
projected indices may be active:
\begin{align}
  \widehat{\phi}_{\pi}
    &\leq \widetilde{\phi}_{\pi},
    \label{eq:amo-lb} \\
  \widetilde{\phi}_{\pi}
    &\leq \widehat{\phi}_{\pi},
    \label{eq:amo-ub}
\end{align}
for each combination of projected indices~$\pi$. Because
$\widehat{\phi}_{\pi}\in\{0,1\}$, both inequalities together force
$\widetilde{\phi}_{\pi}\in\{0,1\}$: either no root binary is active
($\widehat{\phi}_{\pi}=0$) or exactly one is
($\widehat{\phi}_{\pi}=1$).

\smallskip\noindent\textit{Example.}
From $XF_{ijmpr}$, suppressing~$m$ and~$r$
($\pi=(i,j,p)$, $\mathcal{S}=\mathcal{M}\times R$)
yields $\widehat{XF}_{ijp}$, which enforces that at most one
mode--RTI-type pair may carry product~$p$ on route $(i,j)$:
\begin{align}
  \widehat{XF}_{ijp}
    &\leq \widetilde{XF}_{ijp}
    && \forall (i,j)\in E,\;p\in P,
    \label{eq:amo-ex-lb} \\
  \widetilde{XF}_{ijp}
    &\leq \widehat{XF}_{ijp}
    && \forall (i,j)\in E,\;p\in P.
    \label{eq:amo-ex-ub}
\end{align}

\end{convention}

%----------------------------------------------------------------------

Table~\ref{tab:notation-examples} lists the derived variables used in
the model. The accent identifies the aggregation rule,
the suppressed-indices column shows which dimensions are collapsed.

\begin{table}[ht]
\centering
\caption{Derived variables used in the model.}
\label{tab:notation-examples}
\small
\begin{tabular}{@{}lllcl@{}}
\toprule
Derived & Root & Suppr.\ & Rule & Interpretation \\
\midrule
$\widetilde{NF}_{ijmr}$
  & $NF_{ijmpr}$ & $p$
  & $\widetilde{\phantom{x}}$
  & Total daily full-RTI flow on $(i,j)$ by mode $m$, type $r$ \\[3pt]
$\widetilde{N}_{i\cdot r}$
  & $N_{ijmr}$ & $j,m$
  & $\widetilde{\phantom{x}}$
  & Daily outflow of RTI type $r$ from node $i$ \\[3pt]
$\widetilde{N}_{\cdot ir}$
  & $N_{jimr}$ & $j,m$
  & $\widetilde{\phantom{x}}$
  & Daily inflow of RTI type $r$ into node $i$ \\[3pt]
$\widecheck{XF}_{ijmr}$
  & $XF_{ijmpr}$ & $p$
  & $\widecheck{\phantom{x}}$
  & Any product uses full RTIs on $(i,j)$ by $m,r$? \\[3pt]
$\widecheck{X}_{ijm}$
  & $X_{ijmr}$ & $r$
  & $\widecheck{\phantom{x}}$
  & Mode $m$ active on $(i,j)$ for any RTI type? \\[3pt]
$\widecheck{X}_{i\cdot}$
  & $X_{ijmr}$ & $j,m,r$
  & $\widecheck{\phantom{x}}$
  & Node $i$ dispatches any flow? (hub activation) \\[3pt]
$\widecheck{XE}_{i\cdot r}$,\;
$\widecheck{XE}_{\cdot ir}$
  & $XE_{ijmr}$ & $j,m$
  & $\widecheck{\phantom{x}}$
  & Node $i$ dispatching\,/\,receiving empties of type $r$? \\[3pt]
$\widehat{XF}_{ijp}$
  & $XF_{ijmpr}$ & $m,r$
  & $\widehat{\phantom{x}}$
  & At most one mode--RTI-type pair for product $p$ on $(i,j)$ \\
\bottomrule
\end{tabular}
\end{table}
